\section{A growing positive window and its Lambert transition}
\label{sec:positive-wing}

The absolute error in Theorem~\ref{thm:all-orders} is uniform in the
lattice location.  Its usefulness is therefore not confined to bounded
standardized displacements.  We now use its explicit kernels to obtain
a relative approximation in growing positive windows.

Throughout this section, put $L=\log n$ and retain the scale $B$, the
parameters $q,\eps$, and the centered coordinate
$\xi=(m-nA/b)/B$ of Theorem~\ref{thm:all-orders}.  Write
\[
 \mathcal P_n(\xi)=B\mathbb P(X_n=m),\qquad
 g(\xi)=\normal(\xi)=\frac{\exp(-\xi^2/2)}{\sqrt{2\pi}}.
\]
Only lattice coordinates are meant when $\mathcal P_n$ occurs.  The
uniform parameter hypotheses of that theorem give
\begin{equation}\label{wing:small-parameters}
 \frac{c}{L}\le q\le\frac{C}{L},\qquad
 |\eps|\le\frac{C}{L},
\end{equation}
with positive constants independent of the parameters.  In particular,
the constants in the statements below are uniform for the scaled order
in a fixed compact real set and for $b$ in a compact subinterval of
$(0,\infty)$.

\subsection{Right tails of the explicit kernels}

We use the notation
\[
 \mathcal H_{2k}^{[j]}(x)
 =\left.\frac{\partial^j}{\partial\alpha^j}
             \mathcal H_\alpha(x)\right|_{\alpha=2k},\qquad
 \mathcal H_\alpha(x)=\frac1{2\pi}
       \int_{\mathbb R}e^{-iyx-y^2/2}(-iy)^\alpha\,dy,
\]
where $\operatorname{Re}\alpha>-1$ and the principal powers are used.

\begin{samepage}
\begin{lemma}[Kernel tails and order derivatives]\label{lem:kernel-wing}
For fixed integers $k\ge1$ and $j\ge1$, as $x\to+\infty$,
\begin{equation}\label{wing:derivative-bound}
 \mathcal H_{2k}^{[j]}(x)
 =O\bigl(x^{-2k-1}(1+\log x)^{j-1}\bigr).
\end{equation}
More precisely, if $a(\alpha)=1/\Gamma(-\alpha)$, then
\begin{equation}\label{wing:leading-polynomial}
 \mathcal H_{2k}^{[j]}(x)
 =x^{-2k-1}\sum_{r=1}^j\binom jr
        a^{(r)}(2k)(-\log x)^{j-r}
       +O\bigl(x^{-2k-11/4}\bigr).
\end{equation}
In particular,
\begin{equation}\label{wing:first-kernel-tail}
 \mathcal H_2^{[1]}(x)
 =-2x^{-3}\left\{1+6x^{-2}+O(x^{-4})\right\}.
\end{equation}
The constants can depend on $k$ and $j$, which are fixed here.
\end{lemma}
\end{samepage}

\begin{proof}
We give a convolution proof so that differentiation with respect to
the order does not rest on an unstated uniformity of a special-function
asymptotic formula.  Fix $k$, set $d=2k+1$, and let
$|\alpha-2k|\le1/4$.  Then
\begin{equation}\label{wing:convolution}
 \mathcal H_\alpha(x)
 =\frac1{\Gamma(d-\alpha)}\int_0^\infty
               t^{d-\alpha-1}g^{(d)}(x-t)\,dt.
\end{equation}
To verify the identity, first insert $e^{-\eta t}$ with $\eta>0$
in the integral before differentiating $d$ times in $x$.  The Fourier
multiplier of the resulting convolution is
$(\eta-iy)^{\alpha-d}$.  Differentiation multiplies it by
$(-iy)^d$.  As $\eta\downarrow0$, this tends to $(-iy)^\alpha$.
It is dominated, up to a constant uniform on the indicated order
disc, by $|y|^{\operatorname{Re}\alpha}e^{-y^2/2}$: here
$\operatorname{Re}(\alpha-d)<0$ and $|\eta-iy|\ge|y|$.
This is integrable.  Dominated convergence on the physical side is
also valid, because the power at $t=0$ has real part greater than
$-1$ and every derivative of $g$ has Gaussian decay at infinity.
This proves \eqref{wing:convolution}, jointly holomorphically in
$\alpha$ on a neighborhood of the closed disc.

Put $p=d-\alpha-1$ and change variables to $u=t-x$.  On
$|u|\le x/2$, expand $(x+u)^p$ in powers of $u/x$ through degree
$d+2M+1$, for any fixed $M\ge0$.  The Taylor remainder is bounded by
\[
 C_{k,M}x^{\operatorname{Re}p-d-2M-2}|u|^{d+2M+2}
\]
uniformly on the order disc.  Integration against
$|g^{(d)}(-u)|$ is harmless.  Replacing the central integrals by
integrals over the whole real line changes the answer by a quantity
bounded by $e^{-x^2/16}$ times a fixed power of $x$.
The part of the original integral where $|u|>x/2$ has the same
type of bound.  For the left part, the integrable factor
$(x+u)^{\operatorname{Re}p}$ near $u=-x$ causes no difficulty;
for the right part, polynomial growth is dominated by the Gaussian.

Integration by parts gives
\[
 \int_{\mathbb R}u^r g^{(d)}(-u)\,du=0
 \quad(r<d\text{ or }r-d\text{ odd}),
\]
and, for $r=d+2\ell$,
\[
 \int_{\mathbb R}u^{d+2\ell}g^{(d)}(-u)\,du
 =\frac{(d+2\ell)!}{2^\ell\ell!}.
\]
It follows that, uniformly on the order disc,
\begin{equation}\label{wing:uniform-tail-expansion}
 \mathcal H_\alpha(x)
 =\sum_{\ell=0}^{M}
       \frac{x^{-\alpha-2\ell-1}}
            {2^\ell\ell!\,\Gamma(-\alpha-2\ell)}
       +O\bigl(x^{-\operatorname{Re}\alpha-2M-3}\bigr).
\end{equation}
The exponentially small terms just described can be retained
separately; they are already absorbed by this remainder.

The remainder in \eqref{wing:uniform-tail-expansion} is
holomorphic in $\alpha$.  Cauchy's formula on the circle of radius
$1/4$ about $2k$ bounds any fixed order derivative of that remainder
by $O(x^{-2k-2M-11/4})$.  With $M=0$, differentiating the leading
term gives \eqref{wing:leading-polynomial}: the $r=0$ term vanishes
because $a(2k)=0$.  This proves \eqref{wing:derivative-bound}.
For \eqref{wing:first-kernel-tail}, use $M=2$ and $k=j=1$.
The reciprocal Gamma function has a simple zero at each
nonpositive integer, and
\[
 \left.\frac{d}{d\alpha}\frac1{\Gamma(-\alpha-2\ell)}
                         \right|_{\alpha=2}
 =-(2+2\ell)!.
\]
The first two terms are therefore $-2x^{-3}$ and $-12x^{-5}$;
the remaining displayed term and the differentiated remainder are
$O(x^{-7})$, as asserted.
\end{proof}

The same tail structure follows from the classical connection formula
for parabolic cylinder functions and their large-variable expansions \cite{dlmf-cylinder}:
$\mathcal H_\alpha(x)=e^{-x^2/4}D_\alpha(-x)/\sqrt{2\pi}$,
\[
 D_\alpha(-x)=\cos(\pi\alpha)D_\alpha(x)
       +\frac{\pi}{\Gamma(-\alpha)}V(-\alpha-1/2,x).
\]
Here the first term contributes Gaussian decay and the second
contributes the power tail.  The convolution proof above supplies
the parameter uniformity and the order derivatives directly.

\subsection{A relative Gaussian and power-tail approximation}

\begin{theorem}[Positive windows of polynomial size in $\log n$]
\label{thm:positive-windows}
Fix $M>0$, and let $r_L\to\infty$ with $r_L\le L^M$.
Uniformly over lattice points and admissible parameters for which
$r_L\le\xi\le L^M$,
\begin{equation}\label{wing:relative-mixture}
 \mathcal P_n(\xi)
 =\left\{g(\xi)+q\xi^{-3}\right\}\{1+o(1)\}.
\end{equation}
Consequently, for each fixed $\delta\in(0,2)$, whenever the stated
intervals are nonempty,
\begin{align}
 \mathcal P_n(\xi)&\sim g(\xi),
 &r_L\le\xi\le\sqrt{(2-\delta)\log L},
                       \label{wing:gaussian-window}\\
 \mathcal P_n(\xi)&\sim q\xi^{-3},
 &\sqrt{(2+\delta)\log L}\le\xi\le L^M.
                       \label{wing:power-window}
\end{align}
No relative conclusion about negative growing displacements is
included in this theorem.
\end{theorem}

\begin{proof}
Choose an integer $N>3M$ in Theorem~\ref{thm:all-orders}.
Its term of total degree one is
$-q\mathcal H_2^{[1]}(\xi)/2$.  By
\eqref{wing:first-kernel-tail}, this equals
$q\xi^{-3}\{1+O(\xi^{-2})\}$.
For every term with total degree $j\ge2$, divide its absolute
value by $q\xi^{-3}$.  Equations
\eqref{wing:small-parameters} and \eqref{wing:derivative-bound}
bound the quotient by a constant times
\[
 q^{k-1}|\eps|^{j-k}
       \xi^{-2k+2}(1+\log\xi)^{j-1}
 \le C_N\left(\frac{1+\log\xi}{L}\right)^{j-1}.
\]
Only finitely many terms occur.  Their sum is $o(1)$ uniformly
because $\log\xi\le M\log L$.
The absolute remainder divided by $q\xi^{-3}$ is
\[
 O\bigl(\xi^3L^{-N}\bigr)=O\bigl(L^{3M-N}\bigr)=o(1).
\]
Finally $\xi^{-2}\le r_L^{-2}\to0$.
We have thus proved the stronger additive statement
$\mathcal P_n(\xi)=g(\xi)+q\xi^{-3}\{1+o(1)\}$,
which implies \eqref{wing:relative-mixture} since both main terms
are positive.

For \eqref{wing:gaussian-window}, the ratio of the power term to
the Gaussian is at most a constant times
$L^{-\delta/2}r_L^{-3}$.  For \eqref{wing:power-window}, the
ratio of the Gaussian to the power term is at most a constant
times $L^{-\delta/2}(\log L)^{3/2}$: the function
$x^3e^{-x^2/2}$ decreases for $x>\sqrt3$.
Both bounds tend to zero, proving the two consequences.
\end{proof}

In the coefficient specialization of Theorem~\ref{thm:all-orders},
$m=n$, and hence the same conclusions apply directly to
$nBR_b^n a_n$.  Coordinates $\xi\le L^M$ force $b\to A$
because $1-A/b=\xi B/n=O(L^{M+1/2}/\sqrt n)$.
Thus these are growing positive windows adjacent to the feedback
threshold.  They are distinct from a fixed supercritical parameter
regime, whose standardized displacement is much larger.

\begin{corollary}[Relative matching with the supercritical coefficient law]
\label{cor:supercritical-matching}
In the diagonal specialization, uniformly for each fixed $M>0$ and
$\delta>0$ in the nonempty range
\[
 \sqrt{(2+\delta)\log L}\le\xi\le L^M,
\]
one has
\begin{equation}\label{eq:supercritical-matching}
 R_b^n a_n\sim
 \frac1{b\{n(1-A/b)\}^{3+\eps}}
 =\frac{b^{s-1}}{(b-A)^s}\,n^{-s}.
\end{equation}
\end{corollary}
\begin{proof}
On this window $E_2(\eps)\to1$ and
$\xi^{-\eps}\to1$ uniformly, since
$|\eps\log\xi|=O(\log L/L)$. Thus the power term in
\eqref{wing:power-window} is relatively equivalent to
$q\xi^{-3-\eps}/E_2(\eps)$. The exact identities
\[
 \frac{q}{E_2(\eps)}=\frac{n}{bB^{2+\eps}},\qquad
 B\xi=n(1-A/b)
\]
give \eqref{eq:supercritical-matching} after division by $nB$.
\end{proof}

There is an exact explanation for the amplitude cancellation.
The first divided-difference correction is
$-q(\HH_{2+\eps}-\HH_2)/(2\eps)$.
Its leading algebraic tail is
$q\xi^{-3-\eps}/E_2(\eps)$ because
$1/\Gamma(-2-\eps)=-2\eps/E_2(\eps)$.
The polynomial-log window makes this term relatively equivalent to
the simpler $q\xi^{-3}$ while preserving the original coefficient
amplitude in \eqref{eq:supercritical-matching}.

\subsection{An exact Lambert location and a two-term profile}

\begin{theorem}[Lambert location of the Gaussian--power transition]
\label{thm:wing-lambert}
For sufficiently small $q>0$, let $x_*(q)>\sqrt3$ be the larger
positive solution of $g(x_*)=q x_*^{-3}$.  It has the exact formula
\begin{equation}\label{wing:lambert-center}
 x_*(q)^2
 =-3W_{-1}\left(-\frac{(2\pi)^{1/3}}3q^{2/3}\right).
\end{equation}
If $Q=\log(1/q)$, then
\begin{equation}\label{wing:expanded-center}
 x_*(q)^2
 =2Q+3\log(2Q)-\log(2\pi)
       +O\left(\frac{\log Q}{Q}\right).
\end{equation}
For the parameters in \eqref{wing:small-parameters}, uniformly
for $\tau$ in compact real intervals and positive lattice coordinates
$\xi_n>0$ satisfying
\begin{equation}\label{wing:transition-coordinate}
 \xi_n^2=x_*(q)^2+2\tau+o(1),
\end{equation}
where the final error is uniform, one has
\begin{equation}\label{wing:transition-profile}
 \frac{\mathcal P_n(\xi_n)}{qx_*(q)^{-3}}
       \longrightarrow 1+e^{-\tau}.
\end{equation}
In particular $x_*(q)\sim\sqrt{2\log\log n}$.
\end{theorem}

\begin{proof}
The function $x^3g(x)$ increases up to $\sqrt3$ and decreases
thereafter, so the stated larger root exists and is unique for
small $q$.  The defining equation is equivalent to
\[
 x_*^2e^{-x_*^2/3}=(2\pi)^{1/3}q^{2/3},
\]
which gives \eqref{wing:lambert-center} with the branch $W_{-1}$ \cite{dlmf-lambert}.
Alternatively, taking logarithms gives the exact scalar equation
\begin{equation}\label{wing:center-equation}
 y=2Q+3\log y-\log(2\pi),\qquad y=x_*^2.
\end{equation}
It first implies $y\sim2Q$.  Substituting
$y_0=2Q+3\log(2Q)-\log(2\pi)$ into its two sides leaves a
residual $O((\log Q)/Q)$, while the derivative of
$y-3\log y$ is $1-3/y$, bounded away from zero at the larger
root.  The mean value theorem proves
\eqref{wing:expanded-center}.

By \eqref{wing:small-parameters}, $Q=\log L+O(1)$.
Therefore \eqref{wing:transition-coordinate} lies in the
window of Theorem~\ref{thm:positive-windows} for any fixed $M>0$
and a suitable $r_L\to\infty$.  Moreover,
\[
 \frac{g(\xi_n)}{qx_*(q)^{-3}}
 =\frac{g(\xi_n)}{g(x_*(q))}
 =e^{-\tau+o(1)},\qquad
 \frac{q\xi_n^{-3}}{qx_*(q)^{-3}}
 =\left(\frac{x_*(q)}{\xi_n}\right)^3\longrightarrow1.
\]
The relative mixture theorem now gives
\eqref{wing:transition-profile}.
\end{proof}

The lattice restriction is immaterial to the transition coordinate:
rounding $nA/b+B\sqrt{x_*(q)^2+2\tau}$ to the nearest
integer changes $\xi_n^2$ by $O(x_*(q)/B)=o(1)$.
For diagonal coefficients, the feedback parameter can instead be
varied on the locally unique branch near $b=A$, choosing the compact
feedback interval to contain $\zeta(2)$ in its interior.  Here is an explicit
verification of the latter assertion.  Hold $n,s$ fixed, put
$V=bB^2/n$, and differentiate the defining scale equation.  It gives
\[
 \beta:=\partial_b\log B
 =-\frac{b^{-1}+2A^2/(b^3V)}{2-q}=O(1),\qquad
 \partial_b\log q=-b^{-1}-(2+\varepsilon)\beta=O(1).
\]
For $\xi(b)=n(1-A/b)/B$, consequently,
\[
 \xi'(b)=\frac{nA}{b^2B}-\xi(b)\beta
                  \asymp\sqrt{n/L}>0
\]
through any window $\xi=O(L^M)$ near $b=A$.  The derivative of
$x_*(q(b))$ is only $O(1/x_*)$, by
\eqref{wing:center-equation}.  Starting from $\xi(A)=0$,
continuity and these derivative bounds yield a unique local
feedback choice for each bounded $\tau$ in
\eqref{wing:transition-coordinate}.

\subsection{Relation to established probability results}

The Gaussian and power-tail competition is an established phenomenon.
Lam, Blanchet, Burch, and Bazant obtain nonclassical corrections to
Gaussian limits for regularly varying densities, express corrections
using Dawson and parabolic cylinder functions, and explain their
connection with heavy-tail large deviations~\cite{lam}.
Berger, Birkner, and Yuan treat large and local large deviations,
including the infinite-variance normal domain with tail index $-2$,
and identify the transition between collective Gaussian behavior and
a single exceptionally large summand~\cite{berger}.

Accordingly, neither the existence of such a competition nor the use
of parabolic cylinder functions is claimed as new here.  The result
established for the present model is the uniform extension through the
varying order $s_n=3+O(1/\log n)$ and varying nonlinear feedback,
with the explicit paired normalization and the finite all-order
Stirling formulas of Theorem~\ref{thm:all-orders}.  The growing windows
above follow from its uniform absolute remainder, and the exact
Lambert coordinate \eqref{wing:lambert-center} describes their
Gaussian--power transition in that normalization.  The fixed-law
results just cited do not, without additional uniform estimates,
supply this triangular-parameter assertion.
